\section{Описание практической части}
\label{sec:Section4} \index{Section4}

\subsection{Интеграция в систему инвентаризации и валидации сетевых инфраструктур}

В рамках работы разработан прототип системы анализа правил фильтрации и объектов политики безопасности --- \emph{RuleAnalyzer}, и был интегрирован в систему инвентаризации и валидации сетевых инфраструктур, разрабатываемую в ООО <<Дипси Нетворкс>>. Эта система собирает данные о сетевых устройствах, отслеживает их изменения и проверяет на соответствие наборов требований; анализ правил фильтрации и объектов политики безопасности --- одна из таких проверок, и анализатор был встроен в систему в качестве отдельного модуля валидации.

Примерная архитектура модулей системы инвентаризации и валидации сетевых инфраструктур приведена на рис.~\ref{fig:pipeline_1};

\begin{figure}[H]
    \centering
    \includesvg[width=0.8\linewidth]{system_architecture.drawio.svg}
    \caption{Схема обработки конфигурации устройства}
    \label{fig:pipeline_1}
\end{figure}

Система состоит из следующих модулей:
\begin{itemize}
    \item \textbf{Web-Server} и \textbf{Facade} --- точки входа пользователя в систему. \tx{Web-Server} отдаёт веб-интерфейс, \tx{Facade} принимает API запросы от пользователя и маршрутизирует их к внутренним модулям, скрывая их устройство от клиента.
    \item \textbf{Devices-Collector} --- модуль сбора данных о сетевых устройствах. Подключается к сетевым устройствам, выполняет на них команды и забирает вывод команд в исходном текстовом формате производителя, затем выполняет разбор и нормализацию выводов, приводя к унифицированной модели сетевого устройства.
    \item \textbf{Devices-API} --- центральный модуль. Принимает собранные нормализованные данные от \tx{Devices-Collector}, сохраняет их в базе данных, обслуживает запросы от \tx{Facade}. При получении нового снимка сети, отправляет полученные данные на анализ в \tx{RuleAnalyzer}, и сохраняет результаты анализа в базу данных.
    \item \textbf{RuleAnalyzer} --- анализатор правил фильтрации и объектов политики безопасности, описанный и реализованный в данной работе (выделен синим цветом).
\end{itemize}

Для интегрирования нового модуля в систему, также потребовались значительные доработки остальных модулей системы --- каждый из них был адаптирован под новый функционал: была расширена модель данных сетевого устройства (см.~\ref{sec:Section3:model}), чтобы она могла содержать данные о правилах фильтрации и объектах политики безопасности. Devices-Collector был адаптирован под особенности сбора правил фильтрации разных платформ производителей. В Devices-API было добавлено сохранение модели политики и отчётов анализа в базе данных, а в Web-Server --- представление отчёта об аномалиях в веб-интерфейсе.

Модуль анализа реализован на языке Go --- том же, на котором написана остальная бэкенд часть системы, что упрощает интеграцию. Ниже подробно рассмотрены сбор и нормализация выводов команд и результаты нагрузочного тестирования \tx{RuleAnalyzer}.

\subsection{Сбор данных с устройств и их нормализация}
\label{sec:Section4:collector}

Анализатор работает не с самим сетевым устройством, а с унифицированной моделью политики безопасности, описанной в~\ref{sec:Section3:model}. Между устройством и моделью находится коллектор --- модуль, который подключается к устройству, забирает его конфигурацию и приводит её к этой модели. Сложность здесь в том, что каждый производитель отдаёт конфигурацию в собственном текстовом формате: вывод команды \texttt{show access-lists} на Cisco IOS, \texttt{display acl all} на Huawei VRP и аналогов на других платформах синтаксически не совпадают между собой. Коллектор скрывает это различие, разбирая вывод каждой платформы её собственными средствами и складывая результат в общую модель.

\paragraph{Этапы обработки.} Конфигурация проходит через конвейер из нескольких этапов.

\begin{enumerate}
    \item \textbf{Сбор сырого вывода.} Коллектор подключается к устройству по SSH (с поддержкой промежуточных jump-узлов) и выполняет набор команд, заданный для платформы. Команды и их состав версионируются: для каждой пары ``операционная система --- версия'' хранится свой список. Результат каждой команды --- сырой текстовый вывод CLI.
    \item \textbf{Разбор вывода шаблоном TextFSM.} Сырой вывод --- это не структурированные данные, а текст, рассчитанный на чтение человеком: таблицы, отступы, многострочные записи. Для первичного разбора применяется TextFSM --- описание конечного автомата, который по шаблону выделяет из текста отдельные записи и их поля. На этом этапе вывод команды превращается в список записей, где каждая запись --- набор полей ``имя --- значение''.
    \item \textbf{Доразбор строковых полей грамматикой.} TextFSM удобно \emph{режет} вывод на записи и табличные поля, но плохо справляется с семантикой \emph{внутри} одной строки. Строка правила ACL --- например, \texttt{30 permit tcp any host 10.0.0.1 eq 443 established} --- остаётся в записи TextFSM единым строковым полем, тогда как для модели из неё нужно извлечь номер правила, действие, протокол, адреса, порты и хвостовые опции по отдельности. Описать все варианты такого синтаксиса одним конечным автоматом TextFSM непрактично, поэтому строковые поля доразбираются отдельным этапом --- грамматическим разбором (см.\ ниже).
    \item \textbf{Сопоставление с моделью.} Разобранные записи отображаются в унифицированную модель политики: формируются правила фильтрации, именованные объекты, точки применения. Здесь же заполняются служебные поля модели --- в частности, \tx{isComplete} и \tx{unknownMatches} (см.~\ref{sec:Section3:model}, ``Полнота нормализации''): если в разобранной строке встретился квалификатор, которому в модели нет соответствия, он не отбрасывается молча, а сохраняется в \tx{unknownMatches}, а флаг \tx{isComplete} правила сбрасывается.
\end{enumerate}

Этапы~2 и~3 вместе и составляют \emph{нормализацию} --- приведение разнородного вывода устройств к единому структурированному представлению, поверх которого уже работает алгоритм анализа.

\paragraph{Грамматический разбор строк правил.} Доразбор строковых полей реализован через формальные грамматики. Для каждого вида разбираемой строки --- строки ACL, описания именованного объекта и т.\,п.\ --- описывается грамматика в виде двух частей:

\begin{itemize}
    \item \emph{лексический анализатор} (сканер) --- разбивает строку на поток токенов: числа, IP-адреса, ключевые слова (\texttt{permit}, \texttt{tcp}, \texttt{eq}, \dots), идентификаторы, знаки пунктуации. Для строки выше поток будет \texttt{tokNUMBER("30") tokACTION("permit") tokTCP("tcp") \dots};
    \item \emph{синтаксический анализатор} (парсер) --- по грамматическим правилам собирает из потока токенов структуру: распознаёт, что за стандартным или расширенным ACL-правилом стоит номер, действие, протокол, пара адресов с масками, порты и набор хвостовых опций, и формирует из них запись с именованными полями.
\end{itemize}

Лексический и синтаксический анализаторы не пишутся вручную, а \emph{генерируются} по формальному описанию грамматики инструментами \tx{golex} и \tx{goyacc} --- реализациями классических генераторов \tx{lex} и \tx{yacc} для языка Go. Грамматика описывается декларативно: для лексера --- регулярными выражениями, сопоставляющими фрагменты текста с типами токенов, для парсера --- продукциями, описывающими допустимые последовательности токенов и действия по сборке результата. Такой подход даёт два преимущества перед разбором вручную. Во-первых, грамматика проверяема: генератор парсера обнаруживает в ней неоднозначности (конфликты вида shift/reduce) ещё на этапе сборки. Во-вторых, поддержка новой платформы или нового диалекта синтаксиса сводится к добавлению отдельной грамматики, не затрагивая разбор остальных платформ.

Грамматики организованы по иерархии ``производитель --- операционная система --- команда'': при разборе конкретного вывода коллектор по платформе устройства и имени команды выбирает соответствующую грамматику. Если для команды требуется несколько разных грамматик, выбор уточняется ещё и по имени применённого TextFSM-шаблона. На момент написания работы реализованы грамматики для разбора списков доступа и групп объектов платформ Cisco (IOS, IOS-XE, IOS-XR, NX-OS) и Huawei (VRP, YunShanOS).

\subsection{Нагрузочное тестирование}
\label{sec:Section4:benchmark}

Чтобы оценить применимость прототипа к конфигурациям реального масштаба, проведено нагрузочное тестирование модуля анализа. Тестирование измеряет, как время работы и потребление памяти зависят от размера набора правил.

\paragraph{Методика.} Для теста синтезировались наборы правил заданного размера --- от $100$ до $100\,000$ правил. Размер набора --- не единственный параметр, влияющий на стоимость анализа: основной вклад в работу вносит построение \emph{связей} между правилами, а их число зависит от того, насколько часто правила пересекаются. Поэтому наборы синтезировались с разной долей аномальных правил --- $0\%$, $3\%$, $5\%$ и $10\%$: при $0\%$ правила не пересекаются и связи почти не строятся, при $10\%$ доля пересекающихся правил высока. Для каждой пары ``размер набора --- доля аномалий'' измерялись время анализа и пиковое потребление памяти.

\begin{figure}[H]
    \centering
    \includegraphics[width=0.8\linewidth]{benchmark_plots_time.png}
    \caption{Время анализа в зависимости от размера набора правил}
    \label{fig:bench_time}
\end{figure}

\begin{figure}[H]
    \centering
    \includegraphics[width=0.8\linewidth]{benchmark_plots_peak.png}
    \caption{Пиковое потребление памяти в зависимости от размера набора правил}
    \label{fig:bench_peak}
\end{figure}

\paragraph{Результаты.} На рис.~\ref{fig:bench_time} приведено время анализа. На наборах до $20\,000$ правил анализ занимает единицы и первые десятки секунд при любой доле аномалий. На наибольшем наборе в $100\,000$ правил время существенно зависит от доли аномалий: от примерно $50$~секунд при $0\%$ до примерно $620$~секунд при $10\%$. Такое расхождение ожидаемо: сам обход правил линеен по их числу, а дополнительное время уходит на построение связей, которых тем больше, чем чаще правила пересекаются.

Пиковое потребление памяти (рис.~\ref{fig:bench_peak}) растёт с числом правил приблизительно линейно и на наборе в $100\,000$ правил составляет от $13$ до $28$~ГБ в зависимости от доли аномалий. Память расходуется на хранение предикатов правил и накопленных множеств, а также на записи об обнаруженных связях.

\begin{figure}[H]
    \centering
    \includegraphics[width=0.8\linewidth]{benchmark_plots_log.png}
    \caption{Время анализа в двойном логарифмическом масштабе с опорными кривыми $O(n)$ и $O(n^2)$}
    \label{fig:bench_log}
\end{figure}

Чтобы оценить асимптотику, то же время анализа представлено в двойном логарифмическом масштабе (рис.~\ref{fig:bench_log}) вместе с опорными кривыми $O(n)$ и $O(n^2)$. Все замеры лежат между этими кривыми и держатся вблизи $O(n)$, не приближаясь к $O(n^2)$. Это подтверждает заявленную в~\ref{sec:Section3:algorithm} оценку: переход от попарного сравнения правил к накоплению совместной обработки трафика снижает сложность анализа набора с $O(n^2)$ до $O(n)$. Отклонение замеров вверх от строгой прямой $O(n)$ --- вклад построения связей: при высокой доле аномалий одно правило может оказаться связанным с многими, и в худшем случае число связей растёт быстрее линейного. На практике, однако, в реальных конфигурациях доля пересекающихся правил невелика, и анализ наборов реалистичного размера укладывается в приемлемое время.

\newpage